\documentclass[a4paper,addpoints,14pt]{exam}
%\documentclass[answers,addpoints,14pt]{exam}

%tamaños fuentes extra: 8,9,14,17,20
\usepackage{extsizes}

\usepackage[spanish]{babel}
\usepackage[utf8]{inputenc}

\usepackage{array,amsmath,fancybox,amsfonts,graphicx}
\usepackage{amsthm,amssymb,wasysym,latexsym,slashed,anysize}
\usepackage{subfigure,calc}

%\usepackage{fancyhdr,float,hyperref}
%\usepackage{lettrine}

%\usepackage{enumitem}

%entorno para química: texto y math mode
%\newcommand*\chem[1]{\textup{#1}}
\newcommand*\chem[1]{\ensuremath{\mathrm{#1}}}


%display style in long fractions
\newcommand\ddfrac[2]{\frac{\displaystyle #1}{\displaystyle #2}}

%tabla generada
\usepackage[table,xcdraw]{xcolor}

%lewis y química
\usepackage{chemformula,elements}
%chemformula, provides \ch{} with the !(<below>)(<formula>) syntax and \chlewis;
%elements, provides \elconf and \writeelconf.
%ex:
%\ch{
%  !(\elconf{Li})( "\chlewis{180.}{Li}" ) +
%  !(\elconf{F})( "\chlewis{0.90:180:270:}{F}" )
%   ->
%  !(\writeelconf{2})( Li+ ) +
%  !(\writeelconf{2,2+6})( "\chlewis{0:90:180:270:}{F}" {}- )
%}

%multicol
\usepackage{multicol}
  

%símbolo bonito para qed
	%\newcommand{\mc}[3]{\multicolumn{#1}{#2}{#3}}
\newcommand{\halmos}{\vspace*{-6mm}
	
					\begin{flushright}
						\rule{1mm}{2.5mm} 
					\end{flushright}
					\vspace*{-4mm}
						
					  }
\renewcommand{\qed}{\halmos}

 
%página blanco
\usepackage{afterpage}
%uso \afterpage{\blankpage}

\newcommand\blankpage{%
    \null
    \thispagestyle{empty}%
    \addtocounter{page}{0}%
    \newpage}


 %caja bonita
  \usepackage[many]{tcolorbox}
\definecolor{titlebg}{RGB}{80,100,72}
\newtcolorbox{learn}{
  breakable,
  enhanced,
  arc=0pt,
  outer arc=0pt,
  colframe=titlebg,
  colback=titlebg!7,
  overlay unbroken and first={
    \node[
      draw=titlebg,
      fill=titlebg,
      rotate=270,
      anchor=north west,
      text=white,
      font=\bfseries
    ]
    at (frame.north west)  
    {DATOS};
  }
}


\newtcolorbox{formulacion}{
  breakable,
  enhanced,
  arc=0pt,
  outer arc=0pt,
  colframe=titlebg,
  colback=titlebg!7,
  overlay unbroken and first={
    \node[
      draw=titlebg,
      fill=titlebg,
      rotate=270,
      anchor=north west,
      text=white,
      font=\bfseries
    ]
    at (frame.north west)  
    {\hspace*{3ex}};
  }
}

\newtcolorbox{problemas}{
  breakable,
  enhanced,
  arc=0pt,
  outer arc=0pt,
  colframe=titlebg,
  colback=titlebg!7,
  overlay unbroken and first={
    \node[
      draw=titlebg,
      fill=titlebg,
      rotate=270,
      anchor=north west,
      text=white,
      font=
    ]
    at (frame.north west)  
    {AVISO};
  }
}
	
\usepackage{anysize}
	%\marginsize{left}{right}{top}{bottom}
    \marginsize{1cm}{.8cm}{1cm}{1cm}
    

\usepackage{sidecap}
\usepackage{mwe,float}

%redefiniciones
\pointpoints{punto}{puntos}
\bonuspointpoints{punto}{puntos}
\hqword{Pregunta} %substitutes text for \Question"
%\vpgword{Pág.}%substitutestext for \Page"
\hpword{Máx.}%substitutes text for \Points"
\hsword{Nota}%substitutes text for \Score"
\htword{Total}%substitutes text for \Total:"

\chqword{Pregunta}
\chpword{Máx.}
\chsword{Nota}
\chtword{Total}
\chbpword{Puntos \textit{bonus}}

\renewcommand{\solutiontitle}{\noindent\textbf{Solución:}\par\noindent}


%fuentes
	%\usepackage[T1]{fontenc}
	%\usepackage{fourier}

	\usepackage[T1]{fontenc}
	\usepackage{librebaskerville}


\usepackage{graphicx,graphics}
\usepackage{color}
	%	\colorgrids %grids con color
	%	\definecolor{GridColor}{gray}{0.8}
		%grid
%			\setlength{\gridsize}{5mm}
%   			 \setlength{\gridlinewidth}{0.1pt}

%fracciones en texto más elaboradas
	\usepackage{xfrac}

%mejora en la exposición de guarismos
	\usepackage{numprint}
 
%caption
\usepackage[font=small,labelfont=bf,
labelsep=endash]{caption}
	


\checkboxchar{$\Box$}
\checkedchar{$\blacksquare$}

\usepackage{everypage,draftwatermark} %Marca de agua
\SetWatermarkText{B1\\dist.}
\SetWatermarkScale{2}
\SetWatermarkLightness{0.9}



%cabecera
\header{\includegraphics[scale=.22]{a.png}}       
      {}
       {$f_2$\textbf{B1-{\small Distancia}}\\\textbf{febrero}2025}       
%pie de página

   %enumera en todas menos en la primera
   %\runningfootrule

   %enumera en todas
   %\footrule
   %\footer{}{Página \thepage\ de \numpages}{}
%   \footer{}
%	{Página \thepage\ de \numpages}
%	{\iflastpage{{\footnotesize \textit{Fin del examen}}}{{\footnotesize \textit{Da la vuelta a la hoja}\ldots}}}


		
\begin{document}


%portada
%formato de preguntas
\qformat{\textbf{Pregunta \thequestion}\quad (\thepoints)\hfill} 
\pagenumbering{gobble}

\fbox{\parbox{0.945\linewidth}{\textbf{Nombre y apellidos}:\vspace*{20mm}}}
%\vspace{1ex}

%\textbf{\hspace*{5mm} Sólo una de las respuestas es correcta.\newline \hspace*{11mm} Las preguntas erróneamente contestadas restarán 0.25 puntos.}
\vspace*{\fill}

\begin{learn}
{$\triangleright\; M(\chem{O})\simeq16\;g\,mol^{-1}$\newline
$\triangleright\; M(\chem{N})\simeq14\;g\,mol^{-1},$\newline
$\triangleright\; M(\chem{H})\simeq1\;g\,mol^{-1},$\newline
$\triangleright\; M(\chem{S})\simeq32\;g\,mol^{-1},$\newline
$\triangleright\; M(\chem{Cu})\simeq64\;g\,mol^{-1},$\newline
$\triangleright\; R\simeq 0.082\;atm\,L\,mol^{-1}\,K^{-1}$}
\end{learn}	



\newpage
\pagenumbering{gobble}

\vspace*{-1.5cm}

\begin{problemas}{\small Recuerda que un examen es una \textit{comunicación}: debes invocar las leyes en que te basas para resolver los ejercicios así como \textit{indicar brevemente los pasos y razonamientos} que sigues para obtener los resultados.}
\end{problemas}
\vspace{1ex}


\begin{questions}
\question[1] En un matraz de $1.0\;L$ se introducen $0.02\;mol$ de oxígeno molecular y $0.03\;mol$ de nitrógeno molecular. Si la temperatura es de $20^\circ C$ y asumimos que ambos gases se comportan idealmente, calcula:
	\begin{parts}
		\part[\half] La presión total de la mezcla.
		\part[\half] La fracción molar y la presión parcial de cada gas.
	\end{parts}



\question[2] Estima la molaridad de una disolución:
	\begin{parts}
		\part[1] de \chem{HNO_3} etiquetada al $16\%$ en peso y densidad $1.09\;g\,cm^{-3}$.
		\part[1] que resulta de mezclar $0.5\;L$ de una botella $2\;M$ de ácido sulfúrico con  $10\;cm^3$ de $\chem{H_2SO_4}$ concentrado al $96\%$ en peso y densidad $1.84\;g\,cm^{-3}$.
	\end{parts}



\question[2]   En condiciones normales de presión y temperatura se hacen reaccionar $60\;g$ de un mineral de cobre de $90\%$ de pureza con $100\;mL$ de una solución $1\,M$ de ácido nítrico. Sabiendo que el rendimiento para la formación de $\chem{NO_2}(g)$ es del $90\%$ y que se forma, además, nitrato de cobre{\small (II)} y agua:
    	\begin{parts}
    		\part[1] Ajusta la reacción e identifica el reactivo limitante.
    		\part[1] Suponiendo que se comporta idealmente, calcula  el volumen de dióxido de nitrógeno gaseoso que se obtiene.
    	\end{parts}
  	



	\question[2] El movimiento de un proyectil desde un cierto {\small SR} colocado en el suelo puede caracterizarse por 	el vector posición dado por
	\[\vec{r}(t)=\big(t,\;10+30t-5t^2\big)\;\;\text{\small (SI)}.\]
		\begin{parts}
			\part[1] Calcula la velocidad media en los primeros $2\,s$ de vuelo.
			\part[1] Obtén las ecuaciones cinemáticas del proyectil.
			
		\end{parts}
   
  \question[3] Una bala de cañón sale de su boca con una velocidad de $100\;ms^{-1}$ a $1\;m$ de altura sobre el suelo y formando un ángulo inicial de $30^\circ$ sobre el suelo.
     \begin{parts}
    \part[1] Especifica un {\small SR} y las ecuaciones cinemáticas vectoriales.
    \part[1] Calcula la altura máxima que alcanza la bala.
    \part[1] Determina el alcance del proyectil.
   \end{parts}
%\question[2] La descomposición térmica del clorato de potasio se produce según la reacción $\chem{KClO_3(s)}\rightarrow \chem{KCl(s)}+\chem{O_2(g)}.$
%	\begin{parts}
%		\part[1] Calcule la entalpía estándar de reacción para la descomposición de un mol de clorato.
%		%-22 kJ/mol de clorato. (ajuste 1-> 1 , 3/2)
%		\part[1] Determine la cantidad de calor, medida a presión constante, que intercambia la reacción con el medio al obtenerse $30\;L$ de oxígeno molecular, medidos en condiciones estándar de presión y temperatura. \textit{¿Se ceden o se absorben?}
%		%se ceden 17,9 kJ al entorno.
%	\end{parts}
%\vspace{1ex}
%	
%	{
%	\textbf{Datos:} $\Delta_rH^0(\chem{KClO_3(s)})\simeq -414\;kJ\,mol^{-1},\hspace{2ex}\Delta_rH^0(\chem{KCl(s)})\simeq -436\;kJ\,mol^{-1}.$
%	}



%%
%%%\afterpage{\blankpage}
%%\newpage
%%\vspace*{-7ex}
%%
%%\begin{center}
%%       \fbox{\footnotesize HOJA DE RESPUESTAS A LOS PROBLEMAS}
%%\end{center}

\end{questions}	


\end{document}






